\documentclass[conference]{IEEEtran}
\IEEEoverridecommandlockouts
\usepackage[numbers]{natbib}
\usepackage{algorithm,algpseudocode}
\usepackage{amsmath}
\usepackage{caption}
\newcommand*{\vv}[1]{\vec{\mkern0mu#1}}
\usepackage{algorithmicx}
\usepackage{verbatim}
\usepackage{adjustbox}
\usepackage{graphicx, color}
\usepackage{epstopdf}
\usepackage{float}
\usepackage{booktabs}
\usepackage{tabularx}
\usepackage{multirow}
\usepackage{url}
\usepackage{hyperref}
\usepackage[utf8]{inputenc}
\usepackage[T1]{fontenc}
\usepackage{array}
\usepackage{ragged2e}

\newcolumntype{H}{>{\iffalse}c<{\fi}@{}}
\newcommand\tab[1][1cm]{\hspace*{#1}}
\DeclareUnicodeCharacter{00A0}{~}
\DeclareUnicodeCharacter{0301}{\'{e}}
\DeclareUnicodeCharacter{2212}{-}

\renewcommand{\algorithmicrequire}{\textbf{Input:}}
\renewcommand{\algorithmicensure}{\textbf{Output:}}

\newcommand{\mar}[1]{\textcolor{black!90!white}{#1}}

\setlength{\marginparwidth}{2cm}
\def\BibTeX{{\rm B\kern-.05em{\sc i\kern-.025em b}\kern-.07em
    T\kern-.1667em\lower.7ex\hbox{E}\kern-.125emX}}

\begin{document}

\title{
\begin{minipage}{\textwidth}
    \centering
    {\LARGE\textbf{\textcolor{black}{\uppercase{Optimile: AI-Driven Last-Mile Delivery Solution}}}}
\end{minipage}
}

\author{
\IEEEauthorblockN{Mohmed Emad$^{1}$, Mariam Sherif$^{2}$, Zeina Mohamed$^{3}$,\\
Maryam EssamElDin$^{4}$, Yahia Ahmed$^{5}$}
\IEEEauthorblockA{
mohamed2203042$^{1}$, mariam2200593$^{2}$,\\
zeina2207267$^{3}$, maryam2204664$^{4}$, yahia2202676$^{5}$\{@miuegypt.edu.eg\}}
}

\maketitle

% ─────────────────────────────────────────────────────────────
\begin{abstract}
We propose \textbf{Optimile}, a three-layer last-mile delivery optimization
framework for emerging-market logistics. A \textit{Prediction Layer} (DBSCAN zone
clustering, HistGradientBoosting travel-time estimation, and geography-aware
synthetic augmentation) produces an ML cost matrix; an \textit{Optimization Layer}
(Bellman-Ford business-logic sequencing, ALNS with adaptive destroy-repair
operators over 400 iterations, and fleet stop-transfer balancing) converts it into
optimized routes; a \textit{Context \& Adaptation Layer} (OSM Overpass + Google
Directions API hazard detection, reactive re-optimization, and XAI audit logging)
handles live operational disruptions. Built on 590 real GPS segments from 33
drivers across Cairo, Alexandria, and Ismailia (2022--2025), augmented with 6,000
geography-aware synthetic segments (6,590 total).
On 590 real GPS segments, Random Forest achieves CV-MAE\,=\,10.5\,$\pm$\,1.0
min and $R^2$\,=\,0.30 — a \textbf{23.5\%} reduction over the na\"{i}ve mean baseline
(MAE 13.73 min). Physics-augmented training improves CV-MAE to 2.28 min
($R^2$\,=\,0.95); the real-GPS figures are treated as the primary, conservative
performance claim. Component ablation quantifies each contribution: ALNS
delivers 34\% route cost reduction over BF-alone on synthetic instances (9.7\%
on 82 real driver routes); the fleet stop-transfer phase reduces fleet cost by
17\% under naive over-assignment; weather/traffic context features are the
dominant ML contributor, raising CV-MAE by 3.94 min when removed.

\textbf{Keywords:} Last-Mile Delivery, Route Optimization, Adaptive Large
Neighborhood Search, Bellman-Ford, HistGradientBoosting, DBSCAN, Vehicle
Routing Problem, Physics-Based Augmentation, Egyptian Logistics
\end{abstract}

% ─────────────────────────────────────────────────────────────
\section{\textbf{\large Introduction}}
\label{sec:intro}

Last-mile delivery in Egyptian cities presents a distinct set of constraints:
dense urban traffic, informal or imprecise address systems, mixed vehicle
fleets (motorcycles, cars, vans), and courier companies too small to maintain
digitized route logs at research-accessible scale. These factors compound the
standard vehicle routing challenge and make data-driven approaches difficult
to apply directly. A critical barrier is \textit{data scarcity}: small to
mid-size Egyptian courier companies rarely digitize timestamped GPS route
logs. Physics-based synthetic augmentation bridges this gap when real samples
are insufficient~\citep{jie_hybrid_2022, kodikara_comparing_2022}; we adopt
this approach explicitly, keeping real and synthetic evaluation separate
throughout.

Prior work has explored travel time prediction using gradient boosting and
DBSCAN, alongside ALNS and simulated annealing for VRP, showing that weather,
road hazards, and traffic context improve routing
performance~\citep{mara_survey_2022, ilhan_improved_2021,
pereira_approach_2024, liu_time-dependent_2020, moradi_prize-collecting_2025,
he_functionality_2024, abdulrashid_explainable_2024, jie_hybrid_2022,
gmira_managing_2021, kodikara_comparing_2022, prachitmutita_comparison_2020,
pahwa_assessing_2023}. However, none target the Egyptian logistics context
or address the data-scarcity/synthetic-augmentation trade-off explicitly.

We present \textbf{Optimile}, a data-scarcity-aware logistics optimization framework
whose primary contribution is not the novelty of any individual component, but
their coordinated integration into three operational layers: a
\textit{Prediction Layer} (DBSCAN zone clustering, HistGradientBoosting travel-time
estimation, and geography-aware synthetic augmentation); an \textit{Optimization
Layer} (Bellman-Ford business-logic sequencing, ALNS iterative refinement, and
fleet stop-transfer balancing); and a \textit{Context \& Adaptation Layer}
(dual-source hazard detection, reactive re-optimization, and XAI audit logging).
The framework is built on 590 real GPS segments from 33 drivers across Cairo,
Alexandria, and Ismailia (2022--2025), augmented with 6,000 geography-aware
synthetic segments (6,590 total).

Each module is lightweight and interpretable, and the framework degrades
gracefully when contextual signals (traffic, weather, API data) are unavailable
— a deliberate choice reflecting the realities of small Egyptian operators.

\section{\textbf{\large Related Work}}
\label{sec:related}

The contributions of earlier researchers in vehicle routing, adaptive
metaheuristics, hazard-aware routing, and machine learning-based prediction
have shaped the development of the Optimile framework.

Windras Mara et al.~\citep{mara_survey_2022} survey Adaptive Large
Neighborhood Search across 252 applications, classifying destroy operators
such as random, worst, and Shaw removal, and repair operators including
greedy and k-regret insertion. They analyse adaptive weight-update mechanisms
that control operator selection across iterations, showing that score-based
decay strategies are most effective in VRP performance. The Optimile ALNS
adopts this structure using multiple destroy and repair operators with
iterative weight updates over 400 iterations.

Ilhan~\citep{ilhan_improved_2021} proposes an improved simulated annealing
approach for the capacitated vehicle routing problem, introducing a temperature
initialization scheme and geometric cooling schedule to control probabilistic
acceptance of worse solutions. The acceptance mechanism allows the algorithm
to escape local optima by probabilistically accepting non-improving solutions
based on temperature. Optimile integrates this SA criterion within its ALNS
framework using a temperature proportional to solution cost and a gradual
cooling factor to guide convergence over iterations.

Pereira et al.~\citep{pereira_approach_2024} propose ALNS for the
Heterogeneous Fixed Fleet Vehicle Routing Problem with Time Windows
(HFVRPTW), focusing on optimizing fleet operations with different vehicle
capacities under time window constraints. Results show it improves 69.6\% of
benchmark instances with a 0.44\% average cost reduction, running $35\times$
faster than existing approaches. Optimile extends this with a fleet-aware
stop-transfer phase that moves stops between vehicles to reduce total
operational cost before final per-vehicle ALNS optimization.

Liu et al.~\citep{liu_time-dependent_2020} addressed the Time-Dependent VRP
with Time Windows in city logistics by modelling road speed using discrete
congestion levels embedded as multipliers in the routing cost function.
Optimile adopts this idea by applying traffic multipliers
(Low:\,$\times$0.9, Normal:\,$\times$1.0, Medium:\,$\times$1.15,
Heavy:\,$\times$1.35, High:\,$\times$1.60, Jam:\,$\times$1.90)
to travel-time estimates during ALNS cost evaluation.

Moradi and Boroujeni~\citep{moradi_prize-collecting_2025} formulated the
Prize-Collecting Electric VRP with Time Windows, where visiting customers
yields negative cost components that influence optimization behavior.
Optimile adopts this using a Bellman-Ford mechanism that handles negative
weights, assigning urgency rewards ($-15\%$ for stops with $\leq$30 min
time-window slack) and cluster-proximity rewards ($-10\%$ for geographically
nearby stops) to prioritize time-sensitive deliveries.

He et al.~\citep{he_functionality_2024} propose a road network assessment
framework combining flood hazard data with network topology to evaluate road
accessibility. Optimile assigns hazard scores to road segments using OSM
data and weather-weighted risk evaluation to support ALNS re-routing.

Abdulrashid et al.~\citep{abdulrashid_explainable_2024} presented an
explainable AI framework for transport logistics, integrating contrastive
and counterfactual explanations to make routing decisions interpretable.
Optimile adopts this approach through an XAI module that generates
plain-English counterfactuals (e.g., ``Without rerouting, this leg runs
$+$33\% over free-flow time''), confidence badges, and causal event logs for
each routing decision.

Jie, Liu, and Sun~\citep{jie_hybrid_2022} proposed a hybrid algorithm
incorporating weather as multiplicative factors affecting inter-stop travel
times. Optimile adopts a similar mechanism by integrating live weather
updates from OpenWeatherMap and applying travel-time multipliers during
rain, storm, or fog conditions within the ALNS cost function.

Gmira et al.~\citep{gmira_managing_2021} present a reactive re-optimization
framework that triggers local re-optimization when significant travel time
deviations occur. Optimile mirrors this by monitoring road hazards during
execution and re-invoking ALNS from the current vehicle position when delay
thresholds are exceeded.

Kodikara et al.~\citep{kodikara_comparing_2022} investigate ensemble ML
models for delivery time prediction under dynamic conditions, showing that
contextual features --- trip distance, direction, weather --- significantly
improve accuracy. This aligns with Optimile's HistGradientBoosting model for
inter-stop travel time prediction.

Prachitmutita et al.~\citep{prachitmutita_comparison_2020} investigate
density-based clustering to improve last-mile delivery in the Capacitated
VRP, showing that pre-routing clustering reduces planning time and total
travel distance. Optimile uses DBSCAN to group delivery locations into zones
used as features in the prediction model.

Pahwa and Jaler~\citep{pahwa_assessing_2023} assess last-mile resilience
under disruptions through rerouting and fleet reallocation strategies. Optimile
adopts this by incorporating a pre-dispatch hazard detection layer that
triggers ALNS re-optimization before execution.

% ─────────────────────────────────────────────────────────────
\mar{\section{\textbf{\large Proposed Methodology}}}
\label{sec:methodology}

\begin{figure}[H]
\includegraphics[width=\columnwidth]{Fig/methodology.png}
\caption{Proposed System Architecture}
\label{fig:arch}
\end{figure}

% ── 3.1 System Architecture ──────────────────────────────────
\mar{\subsection{System Architecture Overview}}

Figure~\ref{fig:arch} shows the three-layer Optimile pipeline. An optimization
request carries delivery stops (GPS coordinates, fragility flags, time windows),
vehicle type, and optional context (traffic level, weather, incidents).

\textbf{Layer I — Prediction}: Stop coordinates are assigned to delivery zones
via DBSCAN; temporal (cyclic hour, day-of-week) and route-context features are
computed per stop pair. HistGradientBoosting predicts inter-stop travel times
for all $N \times N$ pairs, forming the ML cost matrix.

\textbf{Layer II — Optimization}: Bellman-Ford computes a business-logic
initial ordering (urgency, cluster proximity, traffic, fragility). ALNS then
refines it over 400 iterations with adaptive destroy-repair operators and
Metropolis acceptance. For multi-vehicle scenarios, a transfer phase
redistributes stops from the highest-cost vehicle before per-vehicle ALNS.

\textbf{Layer III — Context \& Adaptation}: OSM Overpass API and Google
Directions API are polled mid-route. When estimated delay exceeds 1.0 minute,
ALNS re-invokes from the vehicle's current position with updated hazard costs.
An XAI module logs plain-English counterfactuals for each routing decision.

Components communicate through a FastAPI backend exposing three endpoints:
\texttt{/optimize} (single vehicle), \texttt{/reoptimize} (mid-route hazard
trigger), and \texttt{/fleet-reoptimize} (multi-vehicle with stop transfer).

% ── Layer I: Prediction ──────────────────────────────────────
\mar{\subsection{Layer I: Travel-Time Prediction}}

\mar{\subsubsection{Datasets \& Augmentation}}

The dataset (\texttt{drivers\_movements.csv}) contains 1,005 GPS
arrival records from 33 drivers across Cairo, Alexandria, and Ismailia
(2022--2025). After filtering malformed timestamps, invalid coordinates, and
gaps outside 1--90 minutes, \textbf{590 valid inter-stop segments} remain.

\textbf{Augmentation Rationale.}A geography-aware synthetic generator bootstrap-resamples real
stop-pairs with $\sigma$=0.001\textdegree{} ($\approx$111 m) perturbation,
enriching each segment with traffic and weather context via a physics formula
calibrated to real GPS statistics (mean 20.3 min, std 18.3 min,
$r_{\text{dist-time}}\!=\!0.546$). This yields 6,000 synthetic segments
preserving the real data's geographic structure, forming a \textbf{6,590-segment}
combined corpus. All ML results are reported separately; real-GPS figures are
the primary claim.

The combined dataset (\texttt{combined\_segments.csv}) contains 6,590 segments,
each described by 16 features covering spatial, temporal, route-context, and
vehicle-type dimensions (Table~1). Real GPS rows lack traffic/weather labels;
these are handled silently by \texttt{OneHotEncoder(handle\_unknown=``ignore'')}.

\begin{table}[H]
\caption{Features of Dataset 2 — Combined Training Segments
(\texttt{combined\_segments.csv})}
\centering
\begin{tabular}{|l|l|p{3.8cm}|}
\hline
\textbf{Feature} & \textbf{Type} & \textbf{Description} \\ \hline
distance\_km         & Numerical   & Haversine distance between stops (km) \\ \hline
hour\_sin            & Numerical   & Sine of cyclic hour-of-day encoding \\ \hline
hour\_cos            & Numerical   & Cosine of cyclic hour-of-day encoding \\ \hline
dow\_sin             & Numerical   & Sine of cyclic day-of-week encoding \\ \hline
dow\_cos             & Numerical   & Cosine of cyclic day-of-week encoding \\ \hline
month                & Numerical   & Month of delivery (1--12) \\ \hline
stop\_index          & Numerical   & Index of current stop in route (0-based) \\ \hline
stops\_in\_route     & Numerical   & Total stops in the driver's route \\ \hline
stops\_remaining     & Numerical   & Stops remaining after current stop \\ \hline
route\_progress      & Numerical   & Ratio of completed legs (0.0--1.0) \\ \hline
is\_same\_cluster    & Binary      & 1 if origin and destination share a zone \\ \hline
bearing\_deg         & Numerical   & Compass direction origin$\to$destination (0--360\textdegree{}) \\ \hline
cluster\_dest\_cat   & Categorical & DBSCAN zone ID of the destination stop \\ \hline
vehicle\_type        & Categorical & Delivery vehicle (car / van / motorcycle) \\ \hline
traffic\_level       & Categorical & Traffic condition (Low/Normal/Medium/High/Jam) \\ \hline
weather              & Categorical & Weather (Sunny/Windy/Sandstorm/Rainy) \\ \hline
time\_to\_next\_stop\_min & Numerical & \textbf{Target}: delivery leg time (min) \\ \hline
\end{tabular}
\end{table}

% ── 3.3 Data Visualization ───────────────────────────────────
\mar{\subsubsection{Data Visualization}}

\paragraph{Real GPS segments.}

The delivery leg time distribution (Figure~\ref{fig:time_dist}) is strongly
right-skewed (skewness 1.57), with the median (14.1 min) well below the mean
(20.3 min). This pattern reflects frequent short urban hops punctuated by
occasional long cross-city legs, and motivates MAE — rather than RMSE — as
the primary evaluation metric to avoid outlier dominance.

\begin{figure}[H]
\centering
\includegraphics[width=0.7\columnwidth]{Fig/delivery_time_dist.png}
\caption{Delivery Time Distribution — right-skewed, mean 20.3 min, median 14.1 min}
\label{fig:time_dist}
\end{figure}

DBSCAN ($\varepsilon\!=\!0.008$\textdegree{}$\!\approx\!0.9$ km, \texttt{min\_samples=3})
identifies 36 delivery zones across all three cities (Figure~\ref{fig:gps_map}).
The $\varepsilon$ value was chosen to match a typical Cairo city block
($\approx$0.9 km), so that stops within the same block cluster together;
\texttt{min\_samples=3} requires at least three GPS stops to form a zone,
preventing singleton clusters from sparse peripheral areas.
29 of these zones appear in at least one training segment as origin or
destination and are used as the \texttt{cluster\_dest\_cat} categorical feature.
The two geographically separated clusters — Alexandria (left) and
Cairo/Ismailia (right) — confirm multi-city driver coverage and illustrate
why zone-aware features improve routing beyond Haversine distance alone.

\begin{figure}[H]
\centering
\includegraphics[width=\columnwidth]{Fig/gps_cluster_map.png}
\caption{GPS Stop Locations Coloured by DBSCAN Delivery Zone (grey = noise)}
\label{fig:gps_map}
\end{figure}

Figure~\ref{fig:by_hour} reveals Cairo's characteristic bimodal rush-hour
pattern: delivery times peak during 8--11h (morning commute) and 15--19h
(evening peak), and are shortest in pre-dawn hours. This empirical pattern
directly informed the piecewise \texttt{hour\_factor} function (0.62--1.35)
used to calibrate synthetic delivery times.

\begin{figure}[H]
\centering
\includegraphics[width=0.7\columnwidth]{Fig/time_by_hour.png}
\caption{Mean Delivery Time by Hour — Cairo rush-hour peaks at 8--11h and 15--19h}
\label{fig:by_hour}
\end{figure}

The scatter plot (Figure~\ref{fig:dist_time}) confirms a moderate distance--time
relationship ($r\!=\!0.546$); the wide residual spread shows that traffic,
weather, and route context are needed to explain the remaining variance.

\begin{figure}[H]
\centering
\includegraphics[width=0.7\columnwidth]{Fig/dist_vs_time.png}
\caption{Distance vs Delivery Time ($r\!=\!0.546$, calibration target confirmed)}
\label{fig:dist_time}
\end{figure}

\paragraph{Combined training dataset.}

Figure~\ref{fig:by_traffic} shows delivery times increasing monotonically from
Low to Jam traffic — validating the physics simulator's multiplicative traffic
factors (Low: 1.40$\times$, Normal: 1.00$\times$, Jam: 0.40$\times$ speed,
resulting in longer times as congestion increases). This monotonic trend is a
necessary condition for model validity.

\begin{figure}[H]
\centering
\includegraphics[width=0.7\columnwidth]{Fig/time_by_traffic.png}
\caption{Delivery Time by Traffic Level — monotonic increase validates simulator multipliers}
\label{fig:by_traffic}
\end{figure}

Weather condition effects follow the expected ordering: Sandstorm produces the
longest delivery times (0.75$\times$ factor) and Sunny conditions the shortest
(1.00$\times$), validating the simulator's weather multipliers alongside the
traffic figure.

\textbf{Synthetic generation} bootstrap-resamples real stop-pairs with
$\sigma$=0.001\textdegree{} coordinate perturbation, preserving the real
distance distribution. Traffic and weather — absent from GPS records — are
sampled from operational priors and used in a calibrated physics formula:
\[
  t = \frac{d}{v_{\text{base}} \times f_{\text{traffic}} \times
  f_{\text{hour}} \times f_{\text{weather}}} \times 60 + \varepsilon
\]
where $v_{\text{base}} \in \{5.8, 8.0, 11.0\}$ km/h (van/car/motorcycle;
effective GPS speed 4.66 km/h including parking and handoff), $f_{\text{traffic}}
\in \{0.40\text{--}1.40\}$, $f_{\text{weather}} \in \{0.75\text{--}1.00\}$,
and $\varepsilon \sim \text{LogNormal}(0, 0.5) - 1$ is right-skewed noise
matching the real GPS shape. Both datasets use 5-fold CV throughout.

\begin{table}[H]
\centering
\caption{Calibration Check — Synthetic vs Real GPS Statistics}
\label{tab:calibration}
\begin{tabular}{|l|c|c|}
\hline
\textbf{Statistic} & \textbf{Real GPS} & \textbf{Synthetic} \\ \hline
Time mean (min)       & 20.3 & 22.2 \\ \hline
Time std (min)        & 18.3 & 23.3 \\ \hline
Distance mean (km)    & 1.578 & 1.586 \\ \hline
Distance std (km)     & 2.138 & 2.000 \\ \hline
Dist--time corr ($r$) & 0.546 & 0.827 \\ \hline
\end{tabular}
\end{table}

Synthetic distance statistics match real GPS within 0.5\% (Table~\ref{tab:calibration}),
confirming geographic structure is preserved. The higher dist--time correlation
in synthetic data reflects the deterministic formula; real GPS shows additional
variance from unobserved factors.

% ── 3.5 Travel Time Prediction Models ────────────────────────
\mar{\subsubsection{Prediction Models}}

Four models were trained and compared using 5-fold CV-MAE, all sharing a
\texttt{StandardScaler} + \texttt{OneHotEncoder(handle\_unknown=``ignore'')}
pipeline. \textbf{ElasticNet} (\texttt{alpha=1.0}, \texttt{l1\_ratio=0.5})
serves as the regularised linear baseline. \textbf{Random Forest}
(\texttt{n\_estimators=300}, \texttt{max\_depth=12}) via bootstrap aggregation
achieves the best real-GPS generalisation (CV-MAE 10.50 min, $R^2$=0.30).
\textbf{Gradient Boosting} (\texttt{n\_estimators=300}, \texttt{lr=0.05},
\texttt{subsample=0.8}) achieves combined-dataset CV-MAE 2.36 min ($R^2$=0.948).
\textbf{HistGradientBoosting} (\texttt{max\_iter=300},
\texttt{max\_depth=6}, \texttt{lr=0.05}) is \textbf{selected as the
production model}: CV-MAE 2.28 min ($R^2$=0.949), with native missing-value
handling (critical for real GPS rows lacking traffic/weather
labels) and faster $N\!\times\!N$ inference.

% ── Layer II: Optimization ───────────────────────────────────
\mar{\subsection{Layer II: Route Optimization}}

This layer converts the ML cost matrix into an optimized stop sequence and,
for multi-vehicle fleets, a load-balanced route assignment. The two components
have \textit{different roles}: Bellman-Ford is a constraint-aware initializer
that encodes urgency, fragility, and proximity rules into the starting visit
order (using negative-weight edges that Dijkstra cannot handle); ALNS then
minimizes route cost from that structured starting point.

\mar{\subsubsection{Stop Sequencing: Bellman-Ford}}
\label{sec:bf}

Bellman-Ford computes the initial stop visit order by encoding business-logic
priority rules as edge weights. Crucially, some rules produce \textit{negative}
weights (urgency rewards, cluster proximity discounts) — making Dijkstra
inapplicable and requiring Bellman-Ford's full relaxation passes. The goal is
\textit{not} route-cost minimization; it is producing a visit order that respects
urgency deadlines and fragility priorities \textit{before} ALNS begins. The weight
$w(i,j)$ between stops $i$ and $j$ is:
\[
  w(i,j) = t_{\text{base}}(i,j) \cdot
  \bigl(1 + \delta_{\text{frag}} + \delta_{\text{traffic}}
  + \delta_{\text{weather}} - \delta_{\text{urgent}}
  - \delta_{\text{cluster}}\bigr)
\]
where $t_{\text{base}}$ is the Haversine-derived travel time when Google
Distance Matrix API data is unavailable, and the modifiers are listed in
Table~\ref{tab:bf_modifiers}. Haversine underestimates road distance in dense
urban grids (straight-line vs.\ actual road path); the Google Distance Matrix
API provides actual road-network travel times when called, and takes precedence.
The algorithm performs $n-1$ relaxation passes
($O(V \cdot E)$) plus one negative-cycle detection pass; detected cycles
fall back to ALNS-only ordering.

\begin{table}[H]
\centering
\caption{Bellman-Ford Edge Weight Modifiers}
\label{tab:bf_modifiers}
\begin{tabular}{|l|l|p{2.5cm}|}
\hline
\textbf{Modifier} & \textbf{Condition} & \textbf{$\delta$} \\ \hline
$\delta_{\text{frag}}$    & Fragile destination             & $+0.20$ \\ \hline
$\delta_{\text{traffic}}$ & Low traffic                     & $-0.05$ \\ \hline
$\delta_{\text{traffic}}$ & Normal traffic                  & $0$     \\ \hline
$\delta_{\text{traffic}}$ & Medium traffic                  & $+0.10$ \\ \hline
$\delta_{\text{traffic}}$ & Heavy traffic                   & $+0.25$ \\ \hline
$\delta_{\text{traffic}}$ & High traffic                    & $+0.40$ \\ \hline
$\delta_{\text{traffic}}$ & Traffic jam                     & $+0.65$ \\ \hline
$\delta_{\text{weather}}$ & Rain / Storm / Fog              & $+0.15$ \\ \hline
$\delta_{\text{urgent}}$  & $\leq$30 min time-window slack  & $-0.15 \times u_f$ \\ \hline
$\delta_{\text{cluster}}$ & Travel time $<$120\,s           & $-0.10$ \\ \hline
\end{tabular}
\end{table}

\mar{\subsubsection{Route Optimization: Adaptive Large Neighborhood Search}}
\label{sec:alns}

\begin{figure}[H]
\centering
\includegraphics[width=\columnwidth]{Fig/alns_flow.png}
\caption{ALNS Destroy-Repair Cycle — operator selection, Metropolis acceptance,
and adaptive weight update repeated over 400 iterations.}
\label{fig:alns_flow}
\end{figure}

The ALNS framework iteratively improves delivery routes through destroy-repair
cycles guided by an adaptive operator weight mechanism and a Metropolis
acceptance criterion (Figure~\ref{fig:alns_flow} and Algorithm~1). Operator
weights are initialized to 1.0 and updated after every iteration via
score-based exponential decay.

\textbf{Destroy Operators:}
\begin{itemize}
\item \textit{Random Removal}: Selects $k\!=\!2$ stops uniformly at random
  and removes them from the current route.
\item \textit{Worst Removal}: Identifies and removes the single stop
  contributing the highest marginal cost to the route.
\item \textit{Fragile-Targeting Removal}: Preferentially targets
  fragile-flagged stops, allowing re-insertion ahead of tighter time windows.
\end{itemize}

\textbf{Repair Operators:}
\begin{itemize}
\item \textit{Greedy Insertion}: Reinserts each removed stop at the position
  yielding the minimum incremental route cost.
\item \textit{Regret Insertion}: Prioritizes stops with the highest regret
  value — the cost difference between the best and second-best insertion
  positions — to avoid suboptimal early placements.
\end{itemize}

\textbf{Acceptance \& Cooling}: A new solution is accepted if it improves the
current best, or with Metropolis probability $P = e^{-\Delta / T}$, where
$\Delta$ is the cost degradation. Initial temperature:
$T_0 = 0.15 \times C_{\text{init}}$. Cooling schedule:
$T \leftarrow 0.995 \times T$ per iteration, converging over 400 iterations.

\textbf{Adaptive Operator Selection}: Each destroy-repair pair maintains a
weight updated as:
\[
  w_{\text{op}} \leftarrow \max\!\left(0.1,\;
  0.8 \cdot w_{\text{op}} + 0.2 \cdot s_{\text{op}}\right)
\]
where $s_{\text{op}} = 5$ for improving solutions, $s_{\text{op}} = 1$ for
accepted non-improving solutions (Metropolis), and $s_{\text{op}} = 0$ for
rejections. Operators are selected via roulette-wheel proportional to their weights.

\textbf{Cost Function}: Total route cost integrates ML-predicted travel times,
traffic and weather multipliers, time-window penalties (wait cost
$0.2 \times t_{\text{wait}}$; late penalty $6.0 \times t_{\text{late}}$),
and fragile surcharge ($2.0 \times t_{\text{travel}}$).

\textbf{Fleet Optimization}: For multi-vehicle scenarios, a transfer phase
moves stops from the highest-cost vehicle to cheaper alternatives before
per-vehicle ALNS, minimizing total fleet ETA.

\begin{algorithm}[H]
\caption{ALNS Route Optimization}
\begin{algorithmic}[1]
\Require stops $S$, vehicle $v$, context $\mathcal{C}$, iters $= 400$
\Ensure optimized route $r^*$
\State $r \leftarrow \text{BellmanFord}(S, \mathcal{C})$;\quad $r^* \leftarrow r$
\State $T \leftarrow 0.15 \times \text{cost}(r)$
\State Initialize operator weights $w_d, w_p \leftarrow 1.0$ for all operators
\For{$i = 1$ \textbf{to} iters}
    \State $d, p \leftarrow \text{RouletteSelect}(w_d, w_p)$
    \State $r' \leftarrow \text{Repair}_p\!\left(\text{Destroy}_d(r,\, k{=}2)\right)$
    \State $\Delta \leftarrow \text{cost}(r') - \text{cost}(r)$
    \If{$\Delta < 0$ \textbf{or} $\text{Uniform}(0,1) < e^{-\Delta / T}$}
        \State $r \leftarrow r'$
        \If{$\text{cost}(r') < \text{cost}(r^*)$}
            \State $r^* \leftarrow r'$
        \EndIf
    \EndIf
    \State $\text{UpdateWeights}(w_d, w_p, \Delta)$
    \State $T \leftarrow 0.995 \times T$
\EndFor
\State \Return $r^*$
\end{algorithmic}
\end{algorithm}

% ── Layer III: Context & Adaptation ─────────────────────────
\mar{\subsection{Layer III: Context \& Adaptation}}

\mar{\subsubsection{Road Hazard Detection}}
\label{sec:hazard}

\begin{figure}[H]
\centering
\includegraphics[width=\columnwidth]{Fig/rerouting_xai.png}
\caption{Live Hazard Detection, Rerouting, and XAI Audit Pipeline —
dual-source hazard inputs, delay threshold trigger, ALNS re-invocation,
and structured XAI record generation.}
\label{fig:rerouting_xai}
\end{figure}

Road hazard detection combines two complementary data sources to compensate
for sparse OSM coverage in Egyptian urban environments.

\textbf{OSM Overpass API}: A bounding-box query is issued for each
optimization request. Road segments are classified into hazard types with
severity scores as shown in Table~\ref{tab:hazard}.

\begin{table}[H]
\centering
\caption{OSM Hazard Classification and Severity Scores}
\label{tab:hazard}
\begin{tabular}{|l|l|c|}
\hline
\textbf{OSM Tag} & \textbf{Hazard Type} & \textbf{Severity} \\ \hline
\texttt{flood\_prone=yes}                    & road\_closed  & 0.85 \\ \hline
\texttt{highway=ford} / \texttt{ford=yes}   & road\_closed  & 0.90 \\ \hline
\texttt{surface=mud}                         & road\_closed  & 0.85 \\ \hline
\texttt{surface=sand}                        & traffic\_jam  & 0.70 \\ \hline
\texttt{surface=dirt} / \texttt{ground}     & traffic\_jam  & 0.55 \\ \hline
\texttt{surface=gravel}                      & traffic\_jam  & 0.45 \\ \hline
\texttt{surface=unpaved}                     & traffic\_jam  & 0.55 \\ \hline
\texttt{surface=grass}                       & traffic\_jam  & 0.50 \\ \hline
\texttt{tunnel=yes}                          & traffic\_jam  & 0.50 \\ \hline
\end{tabular}
\end{table}

\textbf{Google Directions API}: An $N \times N$ pairwise duration matrix is
fetched in batches of 25 origins (Google API limit). Responses use
\texttt{duration\_in\_traffic} when available, falling back to
\texttt{duration}. Responses are cached in 5-minute buckets to minimize API
calls across closely spaced requests.

\textbf{Integration into ALNS}: Detected OSM hazards are merged with any
user-reported incidents. The most severe incident type is injected into the
ALNS cost function with the following delay penalties:
\begin{itemize}
\item \textit{road\_closed}: $+200$ minutes (near-infinite barrier)
\item \textit{accident}: $+60 \times \text{severity}$ minutes
\item \textit{traffic\_jam}: $+35 \times \text{severity}$ minutes
\end{itemize}

\textbf{Re-Optimization Trigger}: If the estimated delay $\geq 1.0$ minute,
ALNS is re-invoked from the vehicle's current GPS position over 300
iterations, with Bellman-Ford providing a refreshed initial ordering under
the updated hazard context.

\textbf{XAI Audit Module}: Every routing decision produces a structured record
with four fields: \textit{counterfactual} (plain-English explanation, e.g.,
``Without rerouting, leg 7$\to$8 runs $+$33\% over free-flow time'');
\textit{confidence} (\textit{HIGH} if incident delay $\geq$50 min or savings
$\geq$15 min, \textit{MEDIUM} if $\geq$15 min or savings $\geq$5 min,
\textit{LOW} otherwise); \textit{event\_log} (per-leg penalty breakdown); and
\textit{savings\_min} (minutes saved vs.\ a reverse-order counterfactual).
Confidence thresholds were calibrated on the 82 real driver routes and align
with the route-size patterns in Table~\ref{tab:ablation_real}.

\mar{\subsection{Performance Metrics}}

Prediction is evaluated using MAE (primary), RMSE, and $R^2$; MAE is chosen
because the right-skewed delivery time distribution (skewness 1.57) makes it
more robust to outliers. Route quality is measured by total route cost and
percentage improvement over baselines. Real GPS and combined datasets are
evaluated in separate CV runs to prevent leakage from physics-calibrated labels.

% ─────────────────────────────────────────────────────────────
\mar{\section{\textbf{\large Results and Analysis}}}
\label{sec:results}

Results are reported separately for 590 real GPS segments and the 6,590-segment
combined dataset.

\subsection{Real GPS Dataset}

\begin{table}[H]
\centering
\caption{CV Results — Real GPS Segments (590 segments, 5-fold)}
\label{tab:real_cv}
\resizebox{\columnwidth}{!}{%
\begin{tabular}{|l|c|c|c|}
\hline
\textbf{Model} & \textbf{CV-MAE (min)} & \textbf{CV-RMSE (min)} & \textbf{CV-$R^2$} \\ \hline
ElasticNet                             & 11.31 & 15.38 & 0.2615 \\ \hline
\textbf{Random Forest}                 & \textbf{10.50} & \textbf{14.85} & \textbf{0.3031} \\ \hline
Gradient Boosting                      & 11.10 & 15.67 & 0.2326 \\ \hline
HistGradientBoosting                   & 11.26 & 15.94 & 0.2019 \\ \hline
\end{tabular}
}
\end{table}

Random Forest achieves CV-MAE 10.50\,$\pm$\,1.0 min and $R^2$\,=\,0.30. The
operational baseline (predict mean) gives MAE 13.73 min; Random Forest's
10.50 min is a \textbf{23.5\% reduction}, providing better ALNS cost-matrix
inputs than a Haversine estimate. ElasticNet's linear structure misses rush-hour
and zone effects, giving the weakest $R^2$ (Figure~\ref{fig:mae_real}).

\begin{figure}[H]
\centering
\includegraphics[width=0.7\columnwidth]{Fig/model_mae_real.png}
\caption{CV-MAE and CV-RMSE Across Models — Real GPS Segments}
\label{fig:mae_real}
\end{figure}

\subsection{Combined Training Dataset}

\begin{table}[H]
\centering
\caption{CV Results — Combined Dataset (6,590 segments, 5-fold)}
\label{tab:comb_cv}
\resizebox{\columnwidth}{!}{%
\begin{tabular}{|l|c|c|c|}
\hline
\textbf{Model} & \textbf{CV-MAE (min)} & \textbf{CV-RMSE (min)} & \textbf{CV-$R^2$} \\ \hline
ElasticNet                             & 10.28 & 14.31 & 0.6084 \\ \hline
Random Forest                          & 3.29  &  6.27 & 0.9243 \\ \hline
Gradient Boosting                      & 2.36  &  5.20 & 0.9477 \\ \hline
\textbf{HistGradientBoosting}          & \textbf{2.28} & \textbf{5.15} & \textbf{0.9486} \\ \hline
\end{tabular}
}
\end{table}

HistGradientBoosting achieves the best combined-dataset CV-MAE (2.28 min,
$R^2$\,=\,0.95), narrowly edging Gradient Boosting (2.36 min, $R^2$\,=\,0.948).
Both tree-boosting models far outperform ElasticNet ($R^2$\,=\,0.61).
\textbf{HistGradientBoosting} is selected as the production model: it
achieves the lowest MAE while providing faster route-matrix inference and
native missing-value handling for real GPS rows lacking traffic/weather labels.

\begin{figure}[H]
\centering
\includegraphics[width=0.7\columnwidth]{Fig/model_r2.png}
\caption{$R^2$ Score Comparison Across Models and Datasets}
\label{fig:r2}
\end{figure}

Figure~\ref{fig:r2} shows tree-model $R^2$ rising from 0.20--0.30 (real GPS)
to 0.92--0.95 (combined), confirming that physics augmentation transfers
delivery time relationships into the training signal.
Figure~\ref{fig:actual_vs_pred} shows good fit for short legs ($<$30 min);
the spread at longer legs reflects unobserved variance (driver behaviour,
unlogged traffic), consistent with real-GPS $R^2$\,=\,0.30.

\begin{figure}[H]
\centering
\includegraphics[width=0.7\columnwidth]{Fig/actual_vs_predicted_real.png}
\caption{Actual vs Predicted Delivery Time — HistGBM, real GPS hold-out (20\%).
Dashed line = perfect prediction.}
\label{fig:actual_vs_pred}
\end{figure}

\subsection{Interpretation and Deployment Trajectory}

The combined $R^2$\,=\,0.95 must be read carefully: 91\% of the corpus is
physics-generated, so the model partly recovers simulator multipliers.
\textbf{Real-GPS figures ($R^2$\,=\,0.30, CV-MAE 10.5 min) are the primary
claim.} The combined model is used in ALNS for its better-calibrated relative
costs under traffic and weather variation.

Table~\ref{tab:learning_curve} quantifies the role of synthetic augmentation:
real GPS alone yields $R^2$\,=\,0.20, rising to 0.95 at 6,590 total. The
fivefold improvement cannot be achieved by collecting more real GPS records
alone. Post-deployment GPS accumulation will enable physics re-calibration
and further improvement.

\begin{table}[H]
\centering
\caption{Augmentation Trajectory — Synthetic Volume vs Model Quality
(590 real records fixed; HistGBM, 5-fold CV).
Generated by \texttt{scripts/ablation\_synthetic\_volume.py}.}
\label{tab:learning_curve}
\begin{tabular}{|c|c|c|c|}
\hline
\textbf{Synthetic added} & \textbf{Total} & \textbf{CV-MAE (min)} & \textbf{CV-$R^2$} \\ \hline
0 (real only)   &   590 & 11.26 & 0.202 \\ \hline
500             & 1,090 &  8.25 & 0.603 \\ \hline
1,000           & 1,590 &  6.24 & 0.740 \\ \hline
2,000           & 2,590 &  4.35 & 0.852 \\ \hline
3,000           & 3,590 &  3.36 & 0.900 \\ \hline
4,500           & 5,090 &  2.68 & 0.931 \\ \hline
6,000 (current) & 6,590 &  2.27 & 0.948 \\ \hline
\end{tabular}
\end{table}

Figure~\ref{fig:aug_traj} shows the same trajectory: monotonic improvement
saturating near 6,000 synthetic records.

\begin{figure}[H]
\centering
\includegraphics[width=0.7\columnwidth]{Fig/augmentation_trajectory.png}
\caption{Augmentation Trajectory — CV-MAE (left axis) and CV-$R^2$ (right axis)
as synthetic volume grows from 0 to 6,000 records (590 real GPS fixed).}
\label{fig:aug_traj}
\end{figure}

% ── 4.4 Component Ablation ───────────────────────────────────
\subsection{Component Ablation Study}

All ablation experiments use fixed random seed 42.

\subsubsection{ML Feature Ablation}

Three configurations evaluated with HistGradientBoosting, 5-fold CV
(Table~\ref{tab:ablation_ml}).

\begin{table}[H]
\centering
\caption{ML Feature Ablation — HistGradientBoosting, 5-fold CV,
Combined Dataset (6,590 segments)}
\label{tab:ablation_ml}
\resizebox{\columnwidth}{!}{%
\begin{tabular}{|l|c|c|c|}
\hline
\textbf{Configuration} & \textbf{CV-MAE (min)} & \textbf{$\pm$ std} & \textbf{CV-$R^2$} \\ \hline
Full system (all features)         & 2.28  & 0.15 & 0.949 \\ \hline
Without weather / traffic features & 6.22  & 0.21 & 0.814 \\ \hline
Distance-only baseline             & 7.41  & 0.33 & 0.748 \\ \hline
\end{tabular}
}
\end{table}

Weather/traffic context features are the dominant contributor: removal raises
CV-MAE by 3.94 min (from 2.28 to 6.22 min) and drops $R^2$ from 0.949 to 0.814.
Distance-only degrades MAE to 7.41 min ($R^2$\,=\,0.748), confirming substantial
ML signal beyond Haversine-only heuristics. DBSCAN zone features are excluded as their impact is largely captured by
\texttt{distance\_km} and \texttt{bearing\_deg}. Note: 91\% physics-generated
data means removing weather/traffic disproportionately hurts synthetic rows;
read alongside real-GPS results (Table~\ref{tab:real_cv}).

\subsubsection{ALNS vs. Bellman-Ford Initialisation Alone}

Table~\ref{tab:ablation_alns} compares greedy nearest-neighbor (NN), BF-only,
and BF\,+\,400 ALNS iterations across 30 instances at four stop-count levels.

\begin{table}[H]
\centering
\caption{NN-only vs.\ BF-only vs.\ BF\,+\,ALNS — Mean Route Cost (30 instances
per stop count, 400 ALNS iterations, \texttt{random\_state=42}).
NN-only column from \texttt{ablation\_alns.py}; BF-only and BF\,+\,ALNS columns
from \texttt{ablation\_bf\_vs\_alns.py}.}
\label{tab:ablation_alns}
\resizebox{\columnwidth}{!}{%
\begin{tabular}{|c|c|c|c|c|c|}
\hline
\textbf{Stops} & \textbf{NN-only (km)} & \textbf{BF-only (km)} & \textbf{BF\,+\,ALNS (km)} & \textbf{Saved vs NN} & \textbf{Saved vs BF} \\ \hline
5  & 3.238 & 3.533 &  3.054 & 5.7\%  & 13.6\% \\ \hline
10 & 5.498 & 6.870 &  4.898 & 10.9\% & 28.7\% \\ \hline
15 & 7.552 & 11.531 & 6.628 & 12.2\% & 42.5\% \\ \hline
20 & 8.690 & 16.091 & 7.880 &  9.3\% & 51.0\% \\ \hline
\multicolumn{4}{|l|}{Mean savings} & \textbf{9.5\%} & \textbf{34.0\%} \\ \hline
\end{tabular}
}
\end{table}

BF\,+\,ALNS outperforms greedy nearest-neighbor dispatch by a mean \textbf{9.5\%}
and BF-alone by a mean \textbf{34.0\%}. The 34\% improvement reflects ALNS's
iterative refinement over 400 destroy-repair cycles, scaling monotonically with
route size. Note that BF-alone records higher cost than NN on instances without
active urgency windows: BF's edge weights are calibrated for business-logic
constraint encoding (urgency, fragile, cluster proximity) rather than
route-cost minimization; ALNS then converts the structured BF ordering into a
cost-optimal sequence. A secondary experiment confirmed that replacing BF initialisation with a
sequential order changes final ALNS cost by only 1.6\%, confirming ALNS
converges robustly from any starting point. BF's contribution is
\textit{structural}: high-urgency stops with $\leq$30 min time-window slack are
visited first and fragile items receive priority reinsertion from iteration~1 —
constraint ordering that ALNS inherits rather than having to discover.

\textbf{Validation on 82 real driver routes} (Cairo, Alexandria, Ismailia;
$\geq$4 stops, sorted by arrival timestamp) confirmed the same pattern on real
GPS coordinates (Table~\ref{tab:ablation_real}).

\begin{table}[H]
\centering
\caption{BF-Only vs. BF\,+\,ALNS on 82 Real Driver Routes (all three cities)}
\label{tab:ablation_real}
\begin{tabular}{|c|c|c|}
\hline
\textbf{Route size} & \textbf{Routes} & \textbf{Mean improvement} \\ \hline
4--6 stops   & 19 &  0.9\% \\ \hline
7--10 stops  & 22 &  4.1\% \\ \hline
11--15 stops & 25 & 13.8\% \\ \hline
16+ stops    & 16 & 21.3\% \\ \hline
\multicolumn{2}{|l|}{Overall (mean / median)} & \textbf{9.7\% / 4.7\%} \\ \hline
\end{tabular}
\end{table}

ALNS gain is minimal for short routes ($<$5\% at 4--10 stops), growing to
21.3\% for 16+ stops. The lower overall mean (9.7\% vs 34\% synthetic) is
expected — real routes skew shorter (median 8 stops).

% ── 4.5 Fleet Transfer Evaluation ────────────────────────────
\subsection{Fleet Stop-Transfer Evaluation}

Table~\ref{tab:ablation_fleet} compares fleet cost with and without the
transfer phase on 30 2-vehicle instances at three imbalance levels (one vehicle
assigned 70--80\% of stops; 300 ALNS iterations).

\begin{table}[H]
\centering
\caption{Fleet Transfer Phase — Total Cost With vs Without Rebalancing}
\label{tab:ablation_fleet}
\begin{tabular}{|c|c|c|c|}
\hline
\textbf{Total stops} & \textbf{Imbalance} & \textbf{No transfer} & \textbf{Improvement} \\ \hline
10 & 70/30 &  6.642 & 22.6\% \\ \hline
16 & 75/25 &  8.615 & 14.2\% \\ \hline
20 & 80/20 & 10.092 & 14.1\% \\ \hline
\multicolumn{3}{|l|}{Mean improvement} & \textbf{17.0\%} \\ \hline
\end{tabular}
\end{table}

Transfer reduces fleet cost by a mean 17.0\%. Figure~\ref{fig:fleet_chart}
shows the cost difference at each imbalance level; the largest gain (22.6\%)
occurs when rebalancing lets both vehicles serve nearby clusters.

\begin{figure}[H]
\centering
\includegraphics[width=0.85\columnwidth]{Fig/fleet_transfer_chart.png}
\caption{Fleet Stop-Transfer — Route Cost With vs Without Rebalancing.
Orange percentages indicate savings; 30 instances per configuration.}
\label{fig:fleet_chart}
\end{figure}

% ── 4.6 System Behavior Illustration ─────────────────────────
\subsection{Illustrative System Behavior}

Figure~\ref{fig:route_comparison} shows a representative 12-stop Cairo route
(real GPS coordinates, driver 64) comparing the unoptimized GPS arrival sequence
against the ALNS-optimized sequence. The optimized order eliminates two
significant cross-city backtracking segments, reducing total route distance by
9.2\% — consistent with the 9.7\% mean improvement measured across 82 real
driver routes in Table~\ref{tab:ablation_real}.

\begin{figure}[H]
\centering
\includegraphics[width=\columnwidth]{Fig/route_comparison.png}
\caption{Before/After Route Optimization — 12-stop real Cairo route.
Left: unoptimized GPS arrival sequence (12.2\,km).
Right: ALNS-optimized sequence (11.1\,km, $-9.2\%$).
Stop numbers indicate visit order; green = start, orange = final stop.}
\label{fig:route_comparison}
\end{figure}

Figure~\ref{fig:rerouting_scenario} shows the full rerouting sequence. An OSM
query returns \texttt{surface=sand} (severity 0.70, classified
\textit{traffic\_jam}), injecting $+24.5$ min ($35\!\times\!0.70$). The delay
exceeds the 1.0 min threshold; ALNS re-invokes over 300 iterations, producing
a rerouted path that avoids the blocked segment. The XAI module outputs:
\textit{counterfactual} ``Without rerouting, leg 7$\to$8 runs $+$33\% over
free-flow time (24.5 min delay). Rerouting saves $\approx$19 min. Confidence:
HIGH''; \textit{event\_log}: base travel 12.1 min + incident penalty 24.5 min;
\textit{savings\_min}: 19.0 min.

\begin{figure}[H]
\centering
\includegraphics[width=\columnwidth]{Fig/rerouting_scenario.png}
\caption{Live Rerouting on a 12-stop Cairo Route.
(a) Optimized pre-dispatch route; (b) OSM hazard detected — \texttt{surface=sand},
$+$24.5 min injected, ALNS re-invoked; (c) Rerouted path, saving $\approx$19 min.}
\label{fig:rerouting_scenario}
\end{figure}

% ── Operational Impact Summary ────────────────────────────────
\subsection{Operational Impact Summary}

\begin{figure}[H]
\centering
\includegraphics[width=\columnwidth]{Fig/operational_impact.png}
\caption{Operational Impact Summary — six key metrics across prediction,
route optimization, fleet balancing, and hazard rerouting.}
\label{fig:op_impact}
\end{figure}

% ── Limitations ───────────────────────────────────────────────
\subsection{Limitations}

Four boundaries should be noted. \textbf{(1) Haversine fallback}: without
Google Distance Matrix API, BF and ALNS use Haversine distance, which
underestimates actual road distance in dense urban grids.
\textbf{(2) No live deployment}: all results are from offline replay and
synthetic simulation — no live dispatch trial, driver study, or SLA
measurement has been conducted; this is the primary gap for future validation.
\textbf{(3) Synthetic data dominance}: the combined $R^2$\,=\,0.95 reflects
a 91\% physics-generated corpus; the model partly recovers simulator
multipliers. Real-GPS $R^2$\,=\,0.30 is the primary external validity claim.
\textbf{(4) No established-solver baseline}: ablation compares against greedy
nearest-neighbor and BF-only, not against OR-Tools or LKH-3.

% ─────────────────────────────────────────────────────────────
\mar{\section{\textbf{\large Conclusion}}}
\label{sec:conclusion}

Optimile demonstrates measurable delivery improvements under the data-scarcity
conditions typical of emerging-market logistics. On 590 real GPS segments,
Random Forest achieves a 23.5\% reduction over the na\"{i}ve mean baseline
(CV-MAE 10.5 min, $R^2$\,=\,0.30). Physics-augmented training ($R^2$\,=\,0.95)
improves ALNS cost-matrix inputs; that metric reflects simulator structure and
is interpreted accordingly.

Component ablation quantifies each contribution: ALNS delivers 34\% route
cost reduction over BF-alone (synthetic), validated at 9.7\% mean on 82 real
driver routes (21.3\% for 16+ stops); fleet stop-transfer reduces fleet cost
by 17\% under naive over-assignment; removing weather/traffic context raises
CV-MAE by 3.94 min, confirming they are the dominant ML contributor.
Real GPS alone yields $R^2$\,=\,0.20; the fivefold improvement to 0.95
requires synthetic augmentation (Table~\ref{tab:learning_curve}). Post-deployment
GPS accumulation will enable physics re-calibration as the corpus grows.

Future work will pursue three directions: (1) a live dispatch trial with
partner couriers, collecting SLA improvement, fuel savings, and driver
feedback to validate the framework against real operational targets; (2) online
retraining as field GPS records accumulate, progressively tightening the
physics calibration; and (3) a dispatcher UI with before/after route maps,
vehicle-balancing visuals, and rerouting overlays, since visual operational
evidence is essential for logistics practitioner adoption.

% ─────────────────────────────────────────────────────────────
\bibliographystyle{IEEEtranN}
\bibliography{references}

\end{document}
